\chapter{Выводы}
\label{ch:conclusions}

Полученные при помощи измерения теплового эффекта Джоуля-Томсона для углекислого газа значения поправок a и b в уравнении Ван-Дер-Ваальса:  
$a  = (1,48 \pm 0,24) \  \frac{H\cdot\text{ м} ^4}{\text{моль} ^2}$, $b = (0.81 \pm 0.19) \cdot 10^{-3} \frac{\text{м} ^3}{\text{моль}}$, и температуры инверсии $T_{\text{инв}} = \frac{2a}{Rb} = (0.44 \pm 0.12) \cdot 10^3 \ K$, разошлись с табличными: $a_{\text{табл}} = 0.365$ $\frac{\text{Н·м}^4}{\text{моль}^2}$, $b_{\text{табл}} = 4.28 \cdot 10^{-3}$ $\frac{\text{м}^3}{\text{моль}}$, $T_{{{\text{инв}}}_{\text{табл}}} = 2 \cdot 10^{3} K$, в несколько раз. Это говорит о том, что уравнение Ван-дер-Ваальса неприменимо для описания поведения углекислого газа.

%с то, не начинать
%значащие цифры

%ван-дер-ваальс 

%стиль результатов:"для того чтобы определить ... мы померили ..." должно быть понятно для чего что мы мерили.

%обсуждение результатов :"построили зависимость, получилась линейная, это свидетельствует о том, что уравнение, которое описывает это поведение правильное.